\Needspace{10\baselineskip}
\section{Application to Silima}

Based on the technical characteristics of Silima and the
established broadcasting practices discussed previously,
we believe that Silima's intended functionality is broadly
consistent with activities permitted under many video game
publishers' broadcasting guidelines.

However, this interpretation is subject to the specific
requirements and restrictions of individual publishers.

\subsection{Characteristics Supporting Our Interpretation}

Several characteristics of Silima support its comparison
to conventional gameplay broadcasting platforms:

\begin{itemize}
    \item \textbf{Centralized Game Execution:}
    Games execute exclusively on the host computer.
    Participants receive gameplay output rather than
    copies of the underlying game software.

    \item \textbf{Broadcasting Functionality:}
    Silima's planned audiovisual streaming functionality
    follows the same general model as conventional
    gameplay broadcasting services, transmitting gameplay
    from a host to connected viewers.

    \item \textbf{Interactive Participation:}
    Silima enables participants to collectively control
    gameplay through voting and shared controller inputs,
    creating an interactive experience beyond passive
    observation.

    \item \textbf{Noncommercial Operation:}
    Silima is intended for noncommercial classroom
    activities. Participants are not charged for
    accessing gameplay.

    \item \textbf{Limited Distribution:}
    Silima operates on a local network with a limited
    audience, rather than broadcasting publicly over
    the internet.

    \item \textbf{No Game Software Distribution:}
    Silima does not provide commercial games, ROMs,
    BIOS files, or other game software to participants.
\end{itemize}

We believe these characteristics establish a reasonable
technical comparison between Silima and established
gameplay broadcasting platforms.

The limited audience and noncommercial nature of Silima
provide additional context, although neither independently
establishes permission to broadcast copyrighted content.

\subsection{Limitations and Unresolved Questions}

Despite these similarities, several considerations
prevent a universal determination that Silima is
authorized under existing publisher policies.

\begin{itemize}
    \item \textbf{Platform Restrictions:}
    Some publishers restrict broadcasting permissions
    to specific websites or streaming services.

    \item \textbf{Individual Versus Institutional Use:}
    Guidelines written for individual content creators
    may not extend to schools or other institutions.

    \item \textbf{Creative Contributions:}
    Some policies require commentary, editing, or other
    original content. Collective gameplay interaction
    may not satisfy these requirements.

    \item \textbf{Third-Party Content:}
    Music and other licensed material appearing within
    games may be subject to separate copyright
    restrictions.

    \item \textbf{Local Broadcasting:}
    Policies addressing public internet broadcasting
    may not explicitly contemplate private network
    transmissions.

    \item \textbf{Game-Specific Restrictions:}
    Individual games may have additional requirements
    that differ from their publisher's general policies.
\end{itemize}

Furthermore, permission to broadcast gameplay does not
automatically authorize copying, downloading, or
circumventing technological protections on game software.

The legality of obtaining and executing games through
emulation is a separate consideration from the
broadcasting of their audiovisual output.

\subsection{Responsibilities of the Host}

Silima provides the technical means for collective
gameplay and its intended broadcasting functionality.
It does not determine whether a particular game may
lawfully be used or broadcast.

The host or institution is responsible for:

\begin{itemize}
    \item Obtaining game software and associated
    resources through lawful means.

    \item Reviewing applicable publisher policies
    and game-specific restrictions.

    \item Ensuring that gameplay presentation complies
    with relevant licensing and copyright requirements.

    \item Avoiding unauthorized distribution of game
    software or protected third-party content.

    \item Complying with applicable institutional
    policies concerning software and media usage.
\end{itemize}

Silima's intended recreational classroom use does not
rely upon an assumption that all activities conducted
within an educational institution qualify as fair use.

Instead, our primary position rests upon the functional
similarity between Silima and established gameplay
broadcasting platforms, evaluated against the express
permissions and limitations established by individual
publishers.

Where those policies do not clearly authorize a
particular use, additional permission or legal
evaluation may be necessary.
